\uuid{0dQK}
\titre{Vrai/faux}
\auteur{bodin ; hanani ; mzari}
\organisation{exo7}
\licence{CC BY-NC-SA 4.0}
\datecreate{2020-01-20}
\niveau{L1}
\module{Algèbre}
\chapitre{Logique, ensemble, raisonnement}
\sousChapitre{Logique}
\difficulte{3}
\theme{Logique, vrai, faux}
\exo7qcmid{lille-1.logique.013}

\contenu{
\begin{question}
\qclassification{100.01}
\qtags{motcle=vrai/faux}
On définit l'assertion "ou exclusif", noté "xou" en disant que "$P$ xou $Q$" est vraie lorsque $P$ est vraie, ou $Q$ est vraie, mais pas lorsque les deux sont vraies en même temps. Quelles sont les assertions vraies ?
\begin{answers}
    \bad{Si "$P$ ou $Q$" est vraie alors "$P$ xou $Q$" aussi.}

    \good{Si "$P$ ou $Q$" est fausse alors "$P$ xou $Q$" aussi.}

    \good{"$P$ xou $Q$" est équivalent à "($P$ ou $Q$) et (non($P$) ou non($Q$))"}

    \bad{"$P$ xou $Q$" est équivalent à "($P$ ou $Q$) ou (non($P$) ou non($Q$))"}
\end{answers}
\begin{explanations}
Commencer par faire la table de vérité de "$P$ ou $Q$".
\end{explanations}
\end{question}
}
